\documentclass[11pt]{amsart}
\usepackage{amsmath,amssymb}
\usepackage{array}
\setlength{\emergencystretch}{3em}
\tolerance=600
\usepackage{tikz-cd}
\usepackage[hidelinks]{hyperref}

\newtheorem{thm}{Theorem}[section]
\newtheorem{lem}[thm]{Lemma}
\newtheorem{cor}[thm]{Corollary}
\newtheorem{prop}[thm]{Proposition}
\theoremstyle{definition}
\newtheorem{defi}[thm]{Definition}
\newtheorem{rem}[thm]{Remark}
\newtheorem{ex}[thm]{Example}
\newtheorem{que}[thm]{Question}

\newcommand{\Same}{\mathrm{Same}}
\newcommand{\Coh}{\mathbf{Coh}}
\newcommand{\Pretop}{\mathbf{Pretop}}
\newcommand{\Sh}{\mathrm{Sh}}
\newcommand{\Mod}{\mathrm{Mod}}
\newcommand{\Set}{\mathbf{Set}}
\newcommand{\Iso}{\mathrm{Iso}}
\newcommand{\FF}{\mathsf F}
\newcommand{\EE}{\mathsf E}
\newcommand{\PP}{\mathsf P}
\newcommand{\cC}{\mathcal C}
\newcommand{\cD}{\mathcal D}
\newcommand{\Fin}{\mathbf{FinSet}}

\begin{document}

\title[Comparing theories in the sameness lattice]{Comparing theories in the lattice of two-out-of-three classes: Morita equivalence as a join of completions, the forced order of coproducts and quotients, and conceptual completeness as an identity of samenesses}
\author{Ryota Shibusawa}
\thanks{Corresponding author. E-mail: \texttt{r-shibusawa@daiichi-koudai.ac.jp}}
\address{Daiichi Institute of Technology}
\email{r-shibusawa@daiichi-koudai.ac.jp}
\keywords{theoretical equivalence, Morita equivalence, definitional equivalence, categorical equivalence, pretopos completion, conceptual completeness, two-out-of-three, localisation, coherent category}
\subjclass[2020]{03G30 (primary), 18F10, 03C07, 18B25 (secondary)}

\begin{abstract}
Definitional equivalence, Morita equivalence (Barrett--Halvorson) and
categorical equivalence of theories are read as classes of
interpretations in the lattice of two-out-of-three classes of the
$2$-category of coherent categories, each completion defining the class
of interpretations it inverts. Morita equivalence is
the class inverted by the pretopos completion, the join of the classes
inverted by the coproduct completion and by the quotient completion,
reached by one alternation with coproducts first; the order is forced,
separated by the propositional theory with one atom, and immaterial for
classical theories. Hence every Morita
equivalence of coherent theories is witnessed by spans of at most four
extensions on each side, three for classical theories with a sort of at
least two elements, and four are needed, by a free theory with two
relations $A\subseteq R$ between two sorts. By the Makkai--Reyes
conceptual completeness theorem the categorical and the Morita classes
coincide at the level of interpretations; hence Hudetz's definable
categorical equivalence implies Morita equivalence when the morphisms
of models are the elementary embeddings, and Morita bi-interpretability
in the sense of Meadows implies Morita equivalence, reversing an arrow
left open in his table. The conceptual distance of
Khaled, Sz\'ekely, Lefever and Friend is the alternation depth of a
join of two classes, finite at every related pair but unbounded.
\end{abstract}

\maketitle

\section{Introduction}

How does one decide that two theories, formulated in different
vocabularies, say the same thing? The question has a long history in
logic (definitional equivalence, synonymy, bi-interpretability) and a
recent one in the philosophy of science, where Barrett and Halvorson
\cite{BH16} proposed \emph{Morita equivalence} as a criterion lying
strictly between definitional equivalence and categorical equivalence,
Tsementzis \cite{Tsem17} identified it with the equivalence of pretopos
completions, Hudetz \cite{Hudetz19} proposed definable categorical
equivalence, McEldowney \cite{McEl20} and Meadows \cite{Meadows24}
mapped the relations between the criteria, and Khaled, Sz\'ekely,
Lefever and Friend \cite{KSLF20} measured the distance between theories
by the number of concepts separating them. The present paper reads
these criteria inside one algebraic object, the \emph{sameness lattice}
$\Same(\cC)$ of the category $\cC$ of theories and interpretations
\cite{Same}: the set of classes of interpretations containing the
isomorphisms, closed under composition and satisfying two-out-of-three,
ordered by inclusion. Each criterion of equivalence is a point of this
lattice, namely the class of interpretations which it declares to be
identifications; comparing criteria is comparing points; and the
lattice operations ask new questions, in particular about joins, which
our earlier work \cite{Fijoin,Cutpoints} has shown to be governed by
the finiteness of an alternation.

\medskip
\noindent\textbf{What is new.} We work in the $2$-category $\Coh$ of
small coherent categories, the syntactic categories of coherent
theories (classical first-order theories enter through their
Morleyisation). A completion $R$ of coherent categories, a reflector
onto a reflective sub-$2$-category, inverts a class $V_R$ of
interpretations, the largest sameness it sends to equivalences
\cite[Prop.~5.1]{Same}.
\begin{enumerate}
\item[(a)] \emph{Morita equivalence is a join of two completions,
reached in one round, in a forced order} (Theorem~\ref{thm:join},
Theorem~\ref{thm:order}). The pretopos completion $\PP$ is the
composite $\EE\FF$ of the completion $\FF$ by finite coproducts and the
completion $\EE$ by quotients of equivalence relations, and $V_\PP=
V_\FF\vee V_\EE$ in $\Same(\Coh)$; $V_\PP$ is Morita equivalence read
as a class of interpretations \cite{Tsem17}. The other order fails:
$\FF\EE\ne\PP$ already for the three-element chain $\bot<p<\top$, the
syntactic category of the coherent propositional theory with one atom
$p$, where the quotient of $\top+\top$ gluing the two copies along $p$
(a sort with two elements which become equal when $p$ holds) is not a
coproduct of quotients. For Boolean coherent categories, i.e.\ for
classical theories, the two completions commute
(Theorem~\ref{thm:boolean}).
\item[(b)] \emph{Lengths of Morita chains} (Theorem~\ref{thm:length}).
With one Morita extension allowed to add all sorts definable over the
current signature, every Morita equivalence of coherent theories is
witnessed by spans of at most four extensions on each side, in the order
products, subsorts, coproducts, quotients, and at most three for
classical theories having a sort with at least two elements, by the
coding of coproduct sorts of \cite[Prop.~5.11]{McEl20}; products must
precede subsorts and coproducts must precede quotients. The bound is
sharp (Theorem~\ref{thm:sharp}): in the free theory with relations
$A\subseteq R\subseteq\sigma\times\tau$ the sort $(R+R)/{\sim_A}$ is not
reached by any three Morita extensions, whatever sorts they add, so that
the theory and the internal theory of its pretopos completion are
connected by no span of chains with three steps on the side of the
theory. The count takes
product and coproduct sorts of any finite arity in one step, as in
\cite{Tsem17}; with the binary sorts of \cite{BH16} no uniform bound
exists (Proposition~\ref{prop:binary}).
\item[(c)] \emph{Categorical equivalence as a class of interpretations
is Morita equivalence} (Theorem~\ref{thm:cc}). By the Makkai--Reyes
conceptual completeness theorem, an interpretation whose induced functor
on categories of models is an equivalence is inverted by the pretopos
completion; this is the observation of Halvorson and Tsementzis
\cite[\S1.3, fn.~8]{HT17} read as the identity $V_{\Mod}=V_\PP$ of two
samenesses. The strict inclusion of Morita equivalence in categorical
equivalence found in \cite[Thm.~5.2]{BH16} therefore concerns pairs of
theories with equivalent model categories but no interpretation inducing
the equivalence; it is not visible in the lattice. As a consequence,
Hudetz's definable categorical equivalence implies Morita equivalence
without the ``external auxiliaries'' of \cite[Thm.~2]{Hudetz19} as soon
as the morphisms of models are the elementary embeddings
(Corollary~\ref{cor:hudetz}), and a Morita interpretation in the sense
of Meadows \cite{Meadows24} whose functor of model categories is an
equivalence is a Morita equivalence; so Morita bi-interpretability
implies Morita equivalence, which reverses an arrow left open in the
table of \cite[\S5.1.1]{Meadows24} (Corollary~\ref{cor:meadows}).
\item[(d)] \emph{Conceptual distance is an alternation depth}
(Proposition~\ref{prop:distance}). The step distance of \cite{KSLF20}
on a cluster network $(X,E,S)$ is the least number of $S$-steps in an
alternating chain, i.e.\ the alternation depth of the join $E\vee S$;
the theorem of \cite{KSLF20} that every finite distance occurs is the
failure of a uniform bound in a join whose alternation is finite at
every pair it relates, the phenomenon studied for closure operators in
\cite{Cutpoints}.
\end{enumerate}
Section~\ref{sec:duality} reads the results as a statement about
duality in the sense of Maruyama's programme \cite{Maruyama16}: a
criterion of equivalence is ``dualisable'' when it is the class inverted
by a completion and the completion is recovered from the models; the
category of models detects the Morita sameness (Makkai--Reyes) and the
ultracategory of models moreover recovers the pretopos (Makkai); the
groupoid of models and isomorphisms, without topology, is left as a
question.

\medskip
\noindent\textbf{What is known.} The hierarchy definitional $\subsetneq$
Morita $\subsetneq$ categorical is \cite{BH16}; Meadows \cite{Meadows24}
arranges twelve criteria in a table with a definitional, a Morita and a
categorical column and rows for isomorphism, equivalence, objective
equivalence and mutual interpretability of model categories, shows that
Morita interpretations induce functors of model categories
\cite[Thm.~42]{Meadows24} and that mixed and pure Morita extensions
give the same equivalence \cite[Cor.~47]{Meadows24}, answers two
questions of \cite{BH16}, and leaves the reversal of most arrows of the
table open; the identification of
Morita equivalence with the equivalence of pretopos completions, and the
explicit span through the syntactic category, products first, then
subsorts, coproducts, quotients, is \cite{Tsem17}; that pretoposes are
the maximal enlargements of a coherent category with the same semantics
is Makkai--Reyes conceptual completeness \cite{MR77}, in the form we use
\cite[Cor.~12, Rem.~13]{LurieNotes}; the ultracategory form is
\cite{Makkai88}. Our contribution is the lattice reading, the join and
order theorems, the counting of extensions, and the consequences for
\cite{BH16,Hudetz19,KSLF20,Meadows24}. The finite-alternation criterion for joins
of classes of closure operators is \cite[Thm.~2.3]{Cutpoints}; the
largest inverted sameness is \cite{Same}. That an equivalence of
model categories induced by an interpretation is a Morita equivalence is
noted in \cite[\S1.3, fn.~8]{HT17} as a consequence of conceptual
completeness \cite[D3.5.9]{Elephant}; Theorem~\ref{thm:cc} is this
remark in the lattice, and its consequences for \cite{Hudetz19} and
\cite{Meadows24} are drawn here.

\section{Samenesses, inverted classes and joins}\label{sec:same}

\begin{defi}[\cite{Same}]
A \emph{sameness} on a category $\cC$ is a class $W$ of morphisms
containing the isomorphisms, closed under composition and satisfying
two-out-of-three: if two of $f$, $g$, $gf$ lie in $W$, so does the
third. $\Same(\cC)$ is the complete lattice of samenesses under
inclusion; meets are intersections. For a functor $F\colon\cC\to\cD$,
$V_F:=\{f:F(f)\text{ invertible}\}$ is the largest sameness $F$ inverts
\cite[Prop.~5.1]{Same}. For a $2$-category we use the same words with
``invertible'' read as ``an equivalence'' and two-out-of-three up to
invertible $2$-cells; $\Same(\cC)$ is then the lattice of such classes.
\end{defi}

Two theories $T$, $T'$ are \emph{$W$-equivalent} if they are connected by
a zig-zag of morphisms in $W$; this is the equivalence relation on
objects induced by $W$. A criterion of equivalence which is not induced
by any sameness is invisible to the lattice; we shall meet one in
Section~\ref{sec:cc}.

The following lemma is the categorical form of the one-round case of
\cite[Thm.~2.3(iv)]{Cutpoints}, where it is shown that for closure
operators $\rho,\rho'$ on a lattice the join of the classes
$W_\rho\vee W_{\rho'}$ is the class of the join $\rho\sqcup\rho'$ iff the
alternation $\rho\rho'\rho\rho'\cdots$ stops at a finite stage at every
point, and that one step suffices iff $\rho\rho'$ is idempotent.

\begin{lem}\label{lem:join}
Let $R_1,R_2$ be reflectors of a $2$-category $\cC$ onto reflective
sub-$2$-categories, with units $\eta_i\colon X\to R_iX$, and suppose
that $R:=R_2R_1$ is again a reflector (its image is reflective with
unit $X\to R_1X\to R_2R_1X$) and that $R$ inverts the units of $R_1$
and of $R_2$. Then $V_R=V_{R_1}\vee V_{R_2}$ in $\Same(\cC)$, and every
$f\in V_R$ is a zig-zag of one $R_1$-unit, one $R_2$-unit, an
equivalence, and the inverses of one $R_2$-unit and one $R_1$-unit.
\end{lem}
\begin{proof}
$V_{R_i}\subseteq V_R$ since $R$ inverts the units of $R_i$: if $R_if$
is an equivalence, the square $R_i f\circ\eta_i^X=\eta_i^Y\circ f$ and
two-out-of-three in $V_R$ give $f\in V_R$. Conversely let $f\colon X\to
Y$ with $Rf$ an equivalence. The unit $X\to RX$ is the composite of
$\eta_1^X\in V_{R_1}$ and $\eta_2^{R_1X}\in V_{R_2}$, so it lies in the
join; likewise for $Y$; and $Rf\circ(X\to RX)=(Y\to RY)\circ f$ up to
an invertible $2$-cell. By two-out-of-three $f$ lies in the join.
\end{proof}

\section{Coherent categories and their completions}\label{sec:coh}

$\Coh$ is the $2$-category of small coherent categories (finite limits,
images, finite unions of subobjects, all stable under pullback) and
coherent functors \cite[A1.4]{Elephant}. A coherent theory $T$ has a
syntactic category $\cC_T$, and interpretations $T\to T'$ are the
coherent functors $\cC_T\to\cC_{T'}$ \cite[D1.4]{Elephant}; every
coherent category is the syntactic category of its internal theory
\cite[\S4]{Tsem17}, so we identify theories with coherent categories.
A coherent category is \emph{positive} if it has finite coproducts which
are disjoint and stable, \emph{effective} if every equivalence relation
is a kernel pair, and a \emph{pretopos} if both. Pretoposes form a
reflective sub-$2$-category $\Pretop\subseteq\Coh$ with reflector
$\PP$, the pretopos completion \cite[8.4]{MR77}, \cite[A3.3]{Elephant}.

We realise the completions inside the sheaf topos. For $\cC\in\Coh$ let
$\Sh(\cC)$ be the sheaves for the coherent topology; the Yoneda
embedding $\cC\to\Sh(\cC)$ is coherent and full. Inside $\Sh(\cC)$ let
\begin{itemize}
\item $\FF\cC$ be the full subcategory of finite coproducts of
representables,
\item $\EE\cC$ be the full subcategory of quotients of representables by
equivalence relations,
\item $\PP\cC$ be the full subcategory of quotients of finite coproducts
of representables by equivalence relations, the coherent objects.
\end{itemize}
$\FF\cC$ is the free positive coherent category on $\cC$, $\EE\cC$ the
free effective one (the exact completion of a regular category
\cite{CV98}), and $\PP\cC$ the pretopos completion; moreover $\EE\FF\cC
=\PP\cC$, i.e.\ \emph{adding coproducts and then quotients gives the
pretopos completion} \cite[8.4]{MR77}, \cite[p.~20]{Tsem17}; that the exact completion
of a positive regular category is positive, so that $\EE\FF\cC$ is a
pretopos, is \cite[Thm.~2.3]{LV01}. All three
are reflectors on $\Coh$ and each inverts the units of the others,
since $\Sh(\FF\cC)\simeq\Sh(\EE\cC)\simeq\Sh(\PP\cC)\simeq\Sh(\cC)$
\cite[D3.3]{Elephant}.

\begin{defi}
Two coherent theories are \emph{Morita equivalent} if they are connected
by a span of Morita extensions in the sense of \cite{BH16,Tsem17}: an
extension by sorts defined as product, coproduct, subsort or quotient
sorts of existing sorts, with the symbols introduced by the definitions,
modulo renaming of symbols. By \cite[Thm.~4.7, Cor.~4.9]{Tsem17} this is
equivalent to the equivalence of pretopos completions, hence to the
$V_\PP$-equivalence of the syntactic categories. We call $V_\PP$ the
\emph{Morita sameness}, $V_\FF$ the \emph{coproduct sameness} and
$V_\EE$ the \emph{quotient sameness}.
\end{defi}

\section{Morita equivalence as a join; the order of coproducts and quotients}\label{sec:join}

\begin{thm}\label{thm:join}
In $\Same(\Coh)$, $V_\PP=V_\FF\vee V_\EE$: the Morita sameness is the
join of the coproduct sameness and the quotient sameness, and the join
is reached by one alternation, coproducts first: every Morita
equivalence $f\colon\cC\to\cD$ is a zig-zag through $\cC\to\FF\cC\to
\EE\FF\cC\xrightarrow{\ \simeq\ }\EE\FF\cD\leftarrow\FF\cD\leftarrow\cD$.
\end{thm}
\begin{proof}
Lemma~\ref{lem:join} with $R_1=\FF$, $R_2=\EE$, $R=\PP=\EE\FF$.
\end{proof}

The lemma applies equally to $R_1=\EE$, $R_2=\FF$ provided $\FF\EE$ is
the pretopos completion. It is not.

\begin{thm}[the order is forced]\label{thm:order}
Let $\cC$ be the three-element chain $\bot<p<\top$, a distributive
lattice, hence a coherent category: the syntactic category of the
coherent propositional theory with one atom $p$ and no axioms. Then
$\EE\cC=\cC$, $\FF\EE\cC=\FF\cC$ is not effective, and $\FF\EE\cC\subsetneq
\PP\cC$. Hence $\FF\EE\ne\PP$, and the alternation of $\EE$ and $\FF$
starting with $\EE$ needs three steps where the alternation starting
with $\FF$ needs two.
\end{thm}
\begin{proof}
A thin category is effective: a reflexive relation on $x$ is a
subobject of $x\times x=x$ containing the identity, so it is $x$, the
kernel pair of the identity; thus $\EE\cC=\cC$. $\Sh(\cC)$ is the topos
of sheaves on the Sierpi\'nski space, equivalently the arrow category
$\Set^{\to}$, an object being a restriction map $X(\top)\to X(p)$; the
representables are $\top=(1\to1)$ and $p=(\emptyset\to1)$, and
$\bot=(\emptyset\to\emptyset)$. Finite coproducts of representables are
the objects $(n\hookrightarrow m)$ with injective restriction, $n\le m$.
The object $\top+\top=(2\to2)$ carries the equivalence relation $R$
identifying the two sections over $p$ and not over $\top$: as a subobject
of $(\top+\top)^2=(4\to4)$ it is $R(\top)=\Delta$, $R(p)=$ everything.
Its quotient in $\Sh(\cC)$ is $(2\to1)$, which is coherent, so it lies
in $\PP\cC$, but its restriction map is not injective, so it is not in
$\FF\cC$. Hence $\FF\EE\cC=\FF\cC\subsetneq\PP\cC$. Moreover $\FF\cC$ is
not effective: otherwise it would be a pretopos containing $\cC$, the
universal property of $\PP\cC$ would give a coherent functor $\PP\cC\to
\FF\cC$ over $\cC$, and composing with the inclusion $\FF\cC\subseteq
\PP\cC$ would give an endofunctor of $\PP\cC$ over $\cC$, necessarily
equivalent to the identity, so that the inclusion would be essentially
surjective, against $(2\to1)\notin\FF\cC$. The last sentence follows: $\EE\FF\cC=\PP\cC$ in two
steps, while $\EE\cC=\cC$, $\FF\EE\cC=\FF\cC$ and only $\EE\FF\EE\cC=\PP\cC$.
\end{proof}

\begin{rem}
The object $(2\to1)$ is a sort with two elements which become equal
when $p$ holds: the quotient of the two-element sort $\top+\top$ by the
relation ``$u\sim v$ iff $u=v$ or $p$''. In a classical theory the same
sort is the coproduct $\{p\}+\{\neg p\}+\{\neg p\}$ of subterminals, a
coproduct of quotients; coherent logic lacks the complement $\neg p$ and
the sort is reached only as a quotient of a coproduct. The same argument
shows that $\FF\EE\ne\PP$ for every distributive lattice with a chain
$\bot<a<x$: in the positive completion the relation $(x,a,a,x)$ on
$(x,x)$ is an equivalence relation whose coequaliser is the fold onto
$(x)$ with kernel pair $(x,x,x,x)$.
\end{rem}

\begin{thm}[Boolean case]\label{thm:boolean}
If $\cC$ is a Boolean coherent category (every subobject has a
complement), then $\FF\EE\cC=\EE\FF\cC=\PP\cC$: every quotient of a finite
coproduct of representables by an equivalence relation of $\FF\cC$, i.e.\
by a coherent equivalence relation in $\Sh(\cC)$, is a finite
coproduct of quotients of representables.
\end{thm}
\begin{proof}
Let $X=X_1+\cdots+X_n$ with $X_i$ representable and $R\subseteq X\times
X$ an equivalence relation in $\FF\cC$, $R=\coprod_{i,j}R_{ij}$ with
$R_{ij}\subseteq X_i\times X_j$ (coproducts are disjoint and stable);
the $R_{ij}$ are subobjects of representables in $\PP\cC$, hence
representable, since the completions add no subobjects to old objects
\cite{CV98}. We
induct on $n$; $n=1$ is the definition of $\EE\cC$. Let $A_1\subseteq
X_1$ be the union of the images of the projections $R_{1j}\to X_1$,
$j\ge2$: a subobject of the representable $X_1$ in $\Sh(\cC)$, hence a
subobject in $\cC$ (the Yoneda embedding is full and preserves images
and finite unions), and let $B_1=X_1\setminus A_1$ be its complement,
again representable. $B_1$ is $R_{11}$-saturated: if $b\,R_{11}\,a$ with
$a\in A_1$, then $a\,R_{1j}\,y$ for some $j\ge2$ and $y$, and transitivity
gives $b\,R_{1j}\,y$, so $b\in A_1$. Hence $R$ restricts to an
equivalence relation $R_{11}|_{B_1}$ on $B_1$ and to an equivalence
relation $R'$ on $X'=A_1+X_2+\cdots+X_n$, with no relations between
$B_1$ and $X'$, and $X/R=B_1/R_{11}+X'/R'$. On $X'$ every element of
$A_1$ is related to an element of $Y:=X_2+\cdots+X_n$, so the composite
$Y\to X'\to X'/R'$ is an epimorphism, every $R'$-class meeting $Y$, hence
a regular epimorphism in the topos $\Sh(\cC)$; its kernel pair is
$R''=R'\cap(Y\times Y)$, the restriction of $R$, so $Y/R''\cong X'/R'$
by effectiveness of equivalence relations in $\Sh(\cC)$. By induction
$X'/R'$ is a
finite coproduct of quotients of representables, and so is $X/R$. All
objects named are representables or their quotients, so the
decomposition takes place in $\FF\EE\cC$.
\end{proof}

\begin{rem}
The induction uses the complement of $A_1$ only; so the conclusion holds
for every coherent $\cC$ in which the images of definable relations
between distinct sorts are complemented. In the chain of
Theorem~\ref{thm:order} the subobject $p\subseteq\top$ is such an image
and has no complement.
\end{rem}

\section{Lengths of Morita chains}\label{sec:length}

We now count extensions. Following \cite[\S2]{Tsem17}, a \emph{Morita
extension} $T\subseteq T^+$ adds to the signature of $T$ new sort
symbols, each defined as the product $S_1\times\cdots\times S_n$ of
finitely many existing sorts, the coproduct $S_1+\cdots+S_n$ of finitely
many existing sorts, a subsort of an existing sort cut out by a formula,
or the quotient of an existing sort by a formula which $T$ proves to be
an equivalence relation, together with the function symbols of the
definitions (projections, injections, inclusion, quotient map), all
definitions being formulas over the \emph{existing} signature; several
sorts may be added at once, but a sort defined from a sort introduced in
the same step is not allowed (\cite[fn.~8]{BH16}: one Morita extension
cannot put predicates on the new sorts). Products and coproducts of
arbitrary finite arity in one step are the convention of
\cite[\S2]{Tsem17}; \cite{BH16} and \cite{Meadows24} define product and
coproduct sorts of two factors only, and the two conventions give the
same Morita equivalence but not the same counts, see
Proposition~\ref{prop:binary}. As in \cite[\S2]{Tsem17}, every theory is
assumed to have a sort which it proves inhabited, which is used to form
a singleton sort; empty definable sets and empty sorts are otherwise
allowed. A \emph{Morita chain of length $k$} is a sequence of $k$ Morita
extensions. Tsementzis's construction \cite[\S4]{Tsem17} gives, for
every coherent theory $T$, a chain $T\subseteq T_1\subseteq\widehat T$
with $T_1$ adding all finite products of sorts and $\widehat T$ adding
all subsorts of these cut out by formulas and all provably functional
relations between them, so that $\cC_{\widehat T}\simeq\cC_T$ and
$\widehat T$ is logically equivalent to the internal theory of $\cC_T$;
and then a chain $T_{\cC}\subseteq T_{\FF\cC}\subseteq T_{\EE\FF\cC}$
adding coproduct sorts and then quotient sorts
\cite[Prop.~4.2]{Tsem17}.

\begin{thm}\label{thm:length}
With product and coproduct sorts of arbitrary finite arity:
\begin{enumerate}
\item[(i)] Two coherent theories are Morita equivalent iff, after
renaming symbols apart, they are connected by a span of Morita chains
of length at most four on each side, with the steps in the order:
products, subsorts, coproducts, quotients.
\item[(ii)] Two of the order constraints cannot be reversed: products
before subsorts, since a subsort of a product (a relation symbol
$R\subseteq\sigma\times\tau$ not a product of unary definable sets) is
not a product of subsorts; and coproducts before quotients, since a
quotient of a coproduct need not be a coproduct of quotients
(Theorem~\ref{thm:order}). Subsorts commute with coproducts (a
subobject of $X+Y$ is $A+B$, by extensivity) and with quotients (a
subsort of $X/E$ is $D/E$ for a saturated $D\subseteq X$, and $D/E$ is a
subsort of $X/(E\cup\Delta_X)$). The bound four is attained, with the
kinds of sorts mixed arbitrarily within a step (Theorem~\ref{thm:sharp}).
\item[(iii)] If $T$ is a classical theory (Boolean syntactic category)
with non-empty sorts and a sort $\sigma$ with $T\vdash\exists u,v\,(u\ne
v)$, then chains of length at most three suffice: products, subsorts,
and then coproducts of quotients in one step.
\end{enumerate}
\end{thm}
\begin{proof}
(i) If $T$ and $T'$ are Morita equivalent their pretopos completions are
equivalent \cite[Cor.~4.9]{Tsem17}. The chain $T\subseteq T_1\subseteq
\widehat T\subseteq T^{(3)}\subseteq T^{(4)}$ of \cite{Tsem17} recalled
above, with $T^{(3)}$ adding the coproduct sorts of the sorts of
$\widehat T$ and $T^{(4)}$ the quotient sorts of the sorts of $T^{(3)}$,
ends in a theory logically equivalent (up to renaming) to the internal
theory of $\PP\cC_T$: by $\EE\FF=\PP$ (Section~\ref{sec:coh}) every
object of $\PP\cC_T$ is a quotient of a finite coproduct of objects of
$\cC_T$, i.e.\ of sorts of $\widehat T$, and every morphism of $\PP\cC_T$
is a provably functional relation between such sorts, definable in
$T^{(4)}$. Two equivalent pretoposes have internal theories which differ
by renaming \cite[p.~20]{Tsem17}. The converse is the definition.

(ii) In the theory with two sorts $\sigma,\tau$ and a relation symbol
$R$, the sort $R$ is the subsort of the product sort $\sigma\times\tau$
cut out by $R(x,y)$; it is not isomorphic to a product of subsorts of
$\sigma$ and $\tau$ in the models where $R$ is not a rectangle. For
coproducts and quotients, Theorem~\ref{thm:order}: in the theory with
one atom $p$, the quotient of $\top+\top$ is not a coproduct of
quotients, and quotient sorts of base sorts give nothing there.

(iii) By Theorem~\ref{thm:boolean} every object of $\PP\cC_T$ is a
finite coproduct $\coprod_i D_i/E_i$ of quotients of definable sets
$D_i$, sorts of $\widehat T$ (reached in two steps). It remains to
obtain such a coproduct as one quotient of one sort of $\widehat T$.
Using $\sigma$ as a tag: $D_1/E_1+D_2/E_2$ is the quotient of the
definable set
\[
D=\{(u,v,\bar x,\bar y):(u=v\wedge\bar x\in D_1)\vee(u\ne v\wedge\bar
y\in D_2)\}
\]
by the equivalence relation identifying $(u,v,\bar x,\bar y)$ with
$(u',v',\bar x',\bar y')$ iff either $u=v$, $u'=v'$ and $\bar x\,E_1\,
\bar x'$, or $u\ne v$, $u'\ne v'$ and $\bar y\,E_2\,\bar y'$. Since
sorts are non-empty the fillers $\bar y$ (resp.\ $\bar x$) exist, and
since $u\ne v$ is satisfiable both tags occur; so the quotient is
$D_1/E_1+D_2/E_2$ in every model, and $D$ is a subsort of a product of
base sorts, a sort of $\widehat T$. For $m$ summands use $m-1$ pairs
$(u_j,v_j)$ of variables of sort $\sigma$ and the mutually exclusive,
jointly exhaustive tags $\theta_i=\bigwedge_{j<i}(u_j\ne v_j)\wedge
(u_i=v_i)$ for $i<m$ and $\theta_m=\bigwedge_{j<m}(u_j\ne v_j)$: let
$D=\{(\bar u,\bar v,\bar x_1,\dots,\bar x_m):\bigvee_i(\theta_i\wedge
\bar x_i\in D_i)\}$ and let $E$ identify two tuples iff they carry the
same tag $\theta_i$ and their $i$-th blocks are $E_i$-related; every tag
is realised and the unused blocks are filled, so $D/E\cong\coprod_i
D_i/E_i$ in every model, in one step and without using the new sorts.
All these are definable over the signature of $\widehat T$, so one more
Morita extension reaches the internal theory of $\PP\cC_T$.
\end{proof}

\begin{rem}
The coding of coproduct sorts by a sort with two provably distinct
elements is not new: McEldowney \cite[Prop.~5.11]{McEl20} uses it to
show that for such theories Morita equivalence implies
bi-interpretability in the model-theoretic sense, and Kamsma
\cite[Fact~8.6]{Kamsma23} reports the fact, as an equivalence, from a
2018 thesis of E.~Washington and transfers it to positive logic; the
two statements are compared after Corollary~\ref{cor:meadows}. What is
added in (iii) is the count: three stages, the coproducts of quotients
being one quotient.
\end{rem}

\begin{rem}
Meadows \cite[Lemma~46, Cor.~47]{Meadows24} proves a normal form of a
different kind: a chain of \emph{mixed} Morita extensions, which may
also add relation and function symbols, is a chain of pure extensions
followed by a single interpretation, the added symbols being postponed
to the end. Theorem~\ref{thm:length} uses pure extensions only and
orders the sort constructors; the two normal forms combine.
\end{rem}

\begin{prop}[binary sorts]\label{prop:binary}
If product and coproduct sorts are restricted to two factors, as in
\cite{BH16,Meadows24}, there is no uniform bound: for every $k$ there
are Morita equivalent theories not connected by any span of Morita
chains of length $k$ on each side.
\end{prop}
\begin{proof}
Let $T$ be the coherent theory with one sort $\sigma$, two constants
$c,d$ and the axiom $c=d\vdash\bot$, and let $T'$ be $T$ with a product
sort $\tau=\sigma^{m}$, $m>2^k$, reached from $T$ by iterated binary
products; $T$ and $T'$ are Morita equivalent. Suppose chains $T\subseteq
T^{(k)}$ and $T'\subseteq T'^{(k)}$ of length $k$ end in definitionally
equivalent theories. Let $M$ be the model of $T$ with $|M_\sigma|=2$ and
$M^{(k)}$ its expansion to $T^{(k)}$: in a binary step the cardinality of
a new sort is at most the square of the largest cardinality of an
existing sort (products $ab\le\max(a,b)^2$, coproducts $a+b\le
\max(a,b)^2$ for $\max(a,b)\ge2$, subsorts and quotients do not
increase cardinality), so every sort of $M^{(k)}$ has at most $2^{2^k}$
elements. Definitionally equivalent theories have the same sorts up to
renaming and a definitional equivalence does not change the domains, so
$M^{(k)}$ is a model $M'^{(k)}$ of $T'^{(k)}$ whose reduct $M'$ is a model
of $T'$; the sort $\sigma$ of $T'$ need not correspond to the sort
$\sigma$ of $T$, but $|M'_\sigma|\ge2$ by the axiom $c=d\vdash\bot$, so
the sort $\tau$ of $M'^{(k)}$ has $|M'_\sigma|^{m}\ge2^{m}>2^{2^k}$
elements, a contradiction.
\end{proof}

\begin{rem}[McEldowney's example]\label{rem:mcel}
McEldowney \cite[p.~20]{McEl20} states, without proof, that the empty
theory in one sort $s$ and its extension $T^\omega$ by binary product
sorts $s^m$ for all $m$ are bi-interpretable in the model-theoretic
sense but not Morita equivalent with binary sorts, ``for the simple
reason that there is no finite sequence of Morita extensions'' from one
to the other. The cardinality argument proves this with the theory $T$
of Proposition~\ref{prop:binary} in place of the empty theory: a span
of length $k$ between $T$ and $T^\omega$ would make the model with
$|M_\sigma|=2$ a model of $T^\omega$ in which every sort has at most
$2^{2^k}$ elements and $\sigma^{m}$ has at least $2^m$, for every $m$.
In the convention of \cite[Def.~2.2]{Tsem17}, where a Morita extension
may add any set of new sorts, $T^\omega$ is a single Morita extension
of $T$ and the two theories are Morita equivalent, as they must be,
having the same pretopos completion. Thus with binary sorts Morita
equivalence is strictly finer than $V_\PP$ once infinitely many sorts
are allowed, and the identification of Morita equivalence with $V_\PP$
(Section~\ref{sec:join}, Theorem~\ref{thm:cc}, \cite[Thm.~4.7]{Tsem17})
is for the convention of \cite{Tsem17}.
\end{rem}

\begin{thm}[the bound four is sharp]\label{thm:sharp}
Let $T_0$ be the coherent theory with sorts $\sigma$, $\tau$ and the unit
sort, relation symbols $A\subseteq\sigma\times\tau$ and
$R\subseteq\sigma\times\tau$, and the single axiom $A(x,y)\vdash R(x,y)$.
Let $\Theta=(R+R)/{\sim_A}=R\cup_A R$ be the object of $\PP\cC_{T_0}$
gluing two copies of $R$ along $A$. Then $\Theta$ is not a sort of any
Morita extension of $T_0$ reached in three steps, whatever kinds of
sorts each step adds. Hence $T_0$ and the internal theory of
$\PP\cC_{T_0}$ are Morita equivalent but connected by no span whose
$T_0$-side has length at most three, the other side being arbitrary; the
chain products, subsorts, coproducts, quotients of
Theorem~\ref{thm:length}(i) is needed.
\end{thm}

A Morita extension $T_k$ of $T_0$ reached in $k$ steps has
$\cC_{T_0}\subseteq\cC_{T_k}$, each new sort being the product,
coproduct, subobject or quotient by which it is defined, so the
inclusion induces an equivalence $\PP\cC_{T_0}\simeq\PP\cC_{T_k}$
(Section~\ref{sec:join}), under which the sorts of $T_k$ are objects of
$\PP=\PP\cC_{T_0}$. Let $S_k$ be the class of objects of $\PP$
isomorphic to a sort of some $k$-step extension; $S_0=\{\sigma,\tau,1\}$
up to isomorphism, and $S_k\subseteq S_{k+1}$. Subobjects and quotients
in $\PP$ of an object of $\cC_{T_0}$ are its definable subobjects and
its quotients by definable equivalence relations \cite{CV98}, the kernel
pair of a quotient being a definable subobject of the square. Call a
\emph{rectangle} a finite product of factors each of which is a sort
(including the unit), a nonzero definable subset of a sort, or a
quotient of a sort by a definable equivalence relation.

\begin{lem}[two-step sorts]\label{lem:twostep}
$S_1$ consists of the finite products of sorts, the definable subsets
of sorts, the finite coproducts of sorts and the quotients of sorts.
Every object of $S_2$ has one of the shapes: ($\alpha$) a subquotient
$D/E$ of a finite product of sorts, $D\subseteq\sigma^a\times\tau^b$
definable and $E$ a definable equivalence relation on $D$, which
includes $0$, the definable sets, the quotients of sorts and of definable
subsets of sorts; ($\beta$) a coproduct of at least one rectangle;
($\gamma$) a quotient $(\sum_a\sigma+\sum_b\tau+\sum_c1)/E$ of a
coproduct of $a+b+c\ge1$ sorts.
\end{lem}
\begin{proof}
A step may add several sorts of different kinds at once, but each is
defined from the existing sorts alone, so the sorts may be considered
one at a time. The description of $S_1$ is then the definition of a
Morita extension, the existing sorts being $\sigma$, $\tau$, $1$. A
second step adds a finite product, a subsort, a finite coproduct or a
quotient of objects of $S_0\cup S_1$. A finite product of such objects distributes over its
coproduct factors and is a finite coproduct of products whose factors
are products of sorts, definable subsets of sorts and quotients of
sorts; each such product is a rectangle, or $0$ if a factor is $0$;
so the product is of shape ($\beta$), or $0$. A subobject of a product
of sorts or of a definable subset is definable, shape ($\alpha$); a
subobject of a coproduct of sorts is, by extensivity, a coproduct of
definable subsets of the summands, shape ($\beta$) or $0$; a subobject
of a quotient $\sigma/E$ is $D/E$ for a saturated definable
$D\subseteq\sigma$, shape ($\alpha$). A coproduct of objects of
$S_0\cup S_1$ is a coproduct of rectangles and of coproducts of sorts,
shape ($\beta$) or $0$. A quotient of a product of sorts or of a
definable subset is of shape ($\alpha$); a quotient of $\sigma/E$ is
$\sigma/E'$ for $E'$ the kernel pair of the composite regular
epimorphism, shape ($\alpha$); a quotient of a coproduct of sorts is of
shape ($\gamma$), or $0$ if there is no summand. Finally $0$, the
empty subobject of a sort, is of shape ($\alpha$).
\end{proof}

\begin{proof}[Proof of Theorem~\ref{thm:sharp}]
Suppose the $T_0$-side of a span has length three: a three-step
extension $T_3$ of $T_0$ is logically equivalent, after renaming, to a
Morita extension $T'$ of the internal theory $T_\PP$ of $\PP$, whose
sorts are all the objects of $\PP$ and whose syntactic category is $\PP$
(Section~\ref{sec:coh}). The logical equivalence identifies $\cC_{T_3}$
with $\cC_{T'}$ and gives an equivalence
$\Phi\colon\PP\simeq\PP\cC_{T'}\simeq\PP\cC_{T_3}\simeq\PP$ under which
every sort of $T_\PP$, that is every object $X$ of $\PP$, is a sort of
$T'$, hence of $T_3$: $\Phi(X)\in S_3$. As $\Phi$ is essentially
surjective and $S_3$ is closed under isomorphism, $S_3$ would contain
every object of $\PP$. We show $\Theta\notin S_3$. Objects of $\PP$ are
evaluated at models of $T_0$, functorially along homomorphisms, the
evaluation being a coherent functor $\PP\to\Set$ which preserves finite
limits, finite coproducts and quotients, so that an epimorphism of
$\PP$ is surjective at every model; coherent formulas are preserved by
homomorphisms. We use the models $F_R=\langle s,t:R(s,t)\rangle$,
$F_A=\langle s,t:A(s,t)\rangle$, the model $M_0$ with
$\sigma=\{s,s'_1,s'_2\}$, $\tau=\{t,t_1,t_2\}$ and
$R=A=\{(s,t_1),(s,t_2),(s'_1,t),(s'_2,t)\}$, the complete bipartite $H_{k,l}$ on
$k+l$ points with $A=\emptyset$, $N=\langle s;t_1,t_2:R(s,t_1),A(s,t_2)\rangle$
and $L'=\langle s_1,s_2;t:R(s_1,t),A(s_2,t)\rangle$, and the one-point
model $1^*$ with $R=A$ everything, which receives a homomorphism from
every model. Then $|\Theta(F_R)|=2$, $|\Theta(F_A)|=1$,
$|\Theta(H_{k,l})|=2kl$, $|\Theta(N)|=|\Theta(L')|=3$. The transpositions
$(t_1\,t_2)$ and $(s'_1\,s'_2)$ are automorphisms of $M_0$ fixing $s$ and
$t$; they generate a group $G$, and no point of $\Theta(M_0)=R(M_0)$ is
fixed by $G$, each edge containing exactly one primed point. Every model
maps to $M_0$ by the homomorphism sending $\sigma$ to $s$ and $\tau$ to
$t_1$, and by the one sending $\sigma$ to $s'_1$ and $\tau$ to $t$.

(a) Every nonzero definable set $D$ (object of $\cC_{T_0}$) is
indecomposable in $\PP$: a complemented subobject of $D$ is a subobject in $\cC_{T_0}$
\cite{CV98}, so a decomposition $D=D'+D''$ is a provable disjointness,
but nonempty $D',D''$ are both nonempty at $1^*$, which has one point
per sort. Hence $R$ is indecomposable, and so is $\Theta$: a complemented
$\Theta'\subseteq\Theta$ meets each copy $R_i$ in $\emptyset$ or $R_i$,
and $R_1\cap R_2=A\ne\emptyset$.

(b) A rectangle $X$ has a $G$-fixed point at $M_0$: take
$s$, $t$, the unit point, or their classes in the quotient factors, and
for a nonzero definable $D\subseteq\sigma$ note $s\in D(M_0)$, a
witness in some model being sent to $s$ by the first homomorphism above;
likewise $t\in D(M_0)$ for $D\subseteq\tau$ by the second, and $U(M_0)=1$
for a nonzero subterminal $U$. A morphism $f\colon X\to\Theta$ of $\PP$
is natural in the model, so $f$ sends $G$-fixed points to $G$-fixed
points; hence $\mathrm{Hom}(X,\Theta)=\emptyset$ for every rectangle $X$,
for every quotient of a rectangle, for every quotient
$(\sum\text{sorts})/E$ of a coproduct of at least one sort (compose
with the cover and restrict to a summand), and for every coproduct of
at least one rectangle (restrict to a summand). The object $0$, the
empty coproduct, is the one exception, $\mathrm{Hom}(0,\Theta)$ being a
point; it is of shape ($\alpha$) and is excluded below by
$|\Theta(F_R)|=2$.

(c) A subquotient of a product of sorts has at most one point at $F_R$;
$|((\Sigma\text{ sorts})/E)(H_{k,l})|\le ak+bl+c$ for $a,b,c$ the numbers
of summands $\sigma,\tau,1$.

By Lemma~\ref{lem:twostep} every object of $S_2$ has shape ($\alpha$),
($\beta$) or ($\gamma$); the products of mixed factors such as
$(\sigma+\sigma)\times D\times D'$ are coproducts of rectangles. A third
step produces $\Theta$ as a quotient, a coproduct, a subsort or a
product of objects of $S_2$, and $0$ plays no part, $\Theta$ being
nonzero.

\emph{Quotient.} An epimorphism $Z\to\Theta$ with $Z\in S_2$: for
($\alpha$) the image has at most one point at $F_R$ by (c); for ($\beta$)
and ($\gamma$) there is no map by (b).
\emph{Coproduct.} $\Theta$ is indecomposable, so $\Theta\cong Z\in S_2$,
and the identity is an epimorphism from $S_2$.
\emph{Subsort.} $\Theta\hookrightarrow Z$: for ($\alpha$) impossible at
$F_R$; for ($\beta$) $\Theta$ embeds in one summand, a rectangle, impossible at $F_R$;
for ($\gamma$) $2kl\le ak+bl+c$ fails for large $k,l$.
\emph{Product.} $\theta\colon\Theta\cong Z_1\times\dots\times Z_n$ with
$Z_i\in S_2$, $n\ge2$ (for $n=1$ see the quotient case). Since
$2=|\Theta(F_R)|=\prod_i|Z_i(F_R)|$, exactly one factor $Z'=Z_j$ has
$|Z'(F_R)|=2$; let $Z=\prod_{i\ne j}Z_i$, so that $\Theta\cong Z\times
Z'$ with projections $p,p'$, $|Z(F_R)|=1$, and $|Z(F_A)|=|Z'(F_A)|=1$
since $|\Theta(F_A)|=1$. Only $Z'\in S_2$ is used below, not $Z\in S_2$.
By (c) $Z'$ is of shape ($\beta$) or ($\gamma$). Shape ($\beta$): two
summands $Z'_1,Z'_2$ of $Z'$ have a point at $F_R$, rectangles having
at most one,
so $Z\times Z'_1$ and $Z\times Z'_2$ are nonempty complementary
subobjects of $\Theta$, against (a). Shape ($\gamma$), $Z'=(\sum_a\sigma
+\sum_b\tau+\sum_c 1)/E$. Suppose $a\ge1$ and use $N$. The homomorphism
$F_A\to N$, $t\mapsto t_2$, puts all copies of $s$, $t_2$ and the unit
points into one $E$-class, since $|Z'(F_A)|=1$; the copies of $t_1$ are
images of elements of $F_R$ under $t\mapsto t_1$, and $|Z'(F_R)|=2$, so
$|Z'(N)|\le3$; as $|\Theta(N)|=3$ is prime and $|Z'(N)|\ge1$, either
$|Z'(N)|=1$ or $|Z(N)|=1$ and $|Z'(N)|=3$. In the second case the two
$F_R$-classes map to two classes disjoint from the class of the
$s$-copies, so neither $F_R$-class contains an $s$-copy, against
$a\ge1$. So $|Z'(N)|=1$. The two copies $r^{(1)},r^{(2)}$ of the edge
$(s,t_1)$ are distinct in $\Theta(N)$, since $A(s,t_1)$ fails; both are
images of the two points of $\Theta(F_R)$ along $F_R\to N$, $t\mapsto
t_1$, so $p(r^{(1)})=p(r^{(2)})$ is the image of the unique point of
$Z(F_R)$, and $p'(r^{(1)})=p'(r^{(2)})$ as $|Z'(N)|=1$; thus $\theta_N$
is not injective. If $a=0$ then $b+c\ge1$ and the same argument runs
in $L'$ with $F_A\to L'$, $s\mapsto s_2$, and $F_R\to L'$, $s\mapsto s_1$:
the copies of $s_2$, $t$ and the unit points form one class, every
$F_R$-class contains a copy of $t$ or of the unit, so $|Z'(L')|=1$, and
the two copies of the edge $(s_1,t)$ collide.
\end{proof}

\begin{rem}
The invariant is point counting and naturality at finitely presented
models of $T_0$; the sharp object is the one anticipated in
Theorem~\ref{thm:order}, a quotient of a coproduct of a non-complemented
definable set, here itself a subsort of a product so that all four kinds
are needed. For classical theories with a two-element sort three steps
suffice (Theorem~\ref{thm:length}(iii)); $T_0$ is not Boolean, $A$ having
no complement in $R$. The four-step chain is the one of \cite[\S4]{Tsem17};
that chains of Morita extensions do not reduce to single extensions is
observed in \cite[\S6.1]{McEl20}, but the least length is not discussed
there or in \cite{BH16}.
\end{rem}

\section{Conceptual completeness: the categorical sameness is the Morita sameness}\label{sec:cc}

Let $\Mod(\cC)$ be the category of coherent functors $\cC\to\Set$ and
natural transformations; for a Boolean $\cC$ (the Morleyisation of a
classical theory) the natural transformations are the elementary
embeddings, the morphisms of models used in \cite{BH16}. $\Mod$ is a
$2$-functor $\Coh^{op}\to\mathbf{CAT}$, and $V_{\Mod}$ is the class of
coherent functors $f$ for which $\Mod(f)$ is an equivalence of
categories. Barrett and Halvorson call two theories \emph{categorically
equivalent} if their categories of models are equivalent, by any
equivalence.

\begin{thm}\label{thm:cc}
$V_{\Mod}=V_\PP$ in $\Same(\Coh)$: an interpretation induces an
equivalence of the categories of models iff it is a Morita equivalence.
\end{thm}
\begin{proof}
The unit $\cC\to\PP\cC$ induces an equivalence $\Mod(\PP\cC)\to\Mod(\cC)$
\cite[Rem.~13]{LurieNotes}, so $\Mod(f)$ is an equivalence iff $\Mod(\PP
f)$ is. If $\Mod(\PP f)$ is an equivalence, then $\PP f$, a morphism of
pretoposes, is an equivalence by the Makkai--Reyes conceptual
completeness theorem \cite[Ch.~7]{MR77}, \cite[Cor.~12]{LurieNotes};
so $f\in V_\PP$. Conversely if $\PP f$ is an equivalence so is $\Mod(\PP
f)$.
\end{proof}

\begin{rem}[the gap of \cite{BH16} is invisible to the lattice]\label{rem:bhgap}
Theorem~5.2 of \cite{BH16} exhibits theories $T_1$ (one sort, countably
many unary predicates $p_i$, the axiom that the sort is a singleton) and
$T_2$ (the same with $q_i$, and the axioms $q_0\to q_i$) whose
categories of models are both discrete with $2^{\aleph_0}$ objects,
hence equivalent, while $T_1$, $T_2$ are not Morita equivalent. By
Theorem~\ref{thm:cc} no interpretation $T_1\to T_2$ or $T_2\to T_1$
induces that equivalence; indeed the interpretation $p_i\mapsto q_i$
induces the inclusion of the models of $T_2$ into those of $T_1$, which
is not essentially surjective. Categorical equivalence as an object
relation is therefore not the equivalence relation induced by any
sameness on $\Coh$ whose restriction to interpretations is $V_{\Mod}$:
the relation induced by $V_{\Mod}$ is Morita equivalence. In the terms
of \cite[Rem.~13]{LurieNotes}, pretoposes are the maximal enlargements
with the same semantics, and the lattice sees exactly the enlargements.
\end{rem}

\begin{cor}[definable categorical equivalence]\label{cor:hudetz}
Let $T_1,T_2$ be classical first-order theories with disjoint
signatures, and take as morphisms of models the elementary embeddings.
If $T_1$ and $T_2$ are definably categorically equivalent in the sense
of \cite[Def.~4]{Hudetz19}, i.e.\ connected by an equivalence of
categories of models given by reconstruction functors, then they are
Morita equivalent.
\end{cor}
\begin{proof}
A reconstruction functor $\Mod(T_1)\to\Mod(T_2)$ is given by a
reconstruction manual \cite[Def.~5--7]{Hudetz19}, which in the
first-order case is an interpretation $\varphi\colon T_2\to T_1$ (each
sort of $T_2$ is sent to a definable subset of a product of sorts of
$T_1$, each symbol to a definable relation, and every model of $T_1$
yields a model of $T_2$, so $T_1$ proves the translated axioms by
completeness), and the functor is $\Mod(\varphi)$ restricted to the
elementary embeddings, which are the natural transformations between
coherent functors on the Morleyisations. If it is an equivalence of
categories, Theorem~\ref{thm:cc} gives $\PP\varphi$ an equivalence, so
the pretopos completions of the Morleyisations are equivalent, and the
theories are Morita equivalent by \cite[Thm.~5.4]{Tsem17}.
\end{proof}

Theorem~2 of \cite{Hudetz19} obtains, under the weaker hypothesis that
only the isomorphisms of models are required to be the
$\Sigma$-isomorphisms, Morita equivalence ``up to external
auxiliaries''; the corollary shows that with elementary embeddings as
morphisms the qualification and the additional definability condition of
\cite[Thm.~3]{Hudetz19} are unnecessary. The groupoid case is the
subject of Question~\ref{que:groupoid}.

Meadows \cite{Meadows24} defines a \emph{Morita interpretation} of $S$
in $T$ as a Morita extension $T^+$ of $T$ together with an
interpretation $t\colon S\to T^+$, shows that it induces a functor
$t^\dagger\colon\mathrm{mod}(T)\to\mathrm{mod}(S)$ on the categories
of models and elementary embeddings, $t^\dagger(M)=t_*(M^+)$ with $M^+$
the canonical expansion of $M$ to $T^+$ \cite[Def.~39, Thm.~42]{Meadows24},
and calls $T$ and $S$ \emph{Morita bi-interpretable} if there are
Morita interpretations in both directions whose induced functors are
quasi-inverse up to definable natural isomorphisms
\cite[Def.~43--44]{Meadows24}. Whether Morita bi-interpretability implies
Morita equivalence is left open in the tables of
\cite[\S5, \S5.1.1]{Meadows24}, which do not use conceptual completeness.

\begin{cor}[Morita interpretations]\label{cor:meadows}
Let $T$, $S$ be classical first-order theories, with elementary
embeddings as morphisms of models. If $T$ Morita interprets $S$ via
$(T^+,t)$ and the induced functor $t^\dagger\colon\mathrm{mod}(T)\to
\mathrm{mod}(S)$ is an equivalence of categories, then $T$ and $S$ are
Morita equivalent. In particular Morita bi-interpretability implies
Morita equivalence, the natural isomorphisms being allowed to be
arbitrary; conversely Morita equivalence is $1$-isomorphism in Meadows's
$2$-category $\mathbf{TH}_{\mathrm{Mor}}$ \cite[Cor.~47]{Meadows24},
hence $1$-equivalence, so the two notions coincide.
\end{cor}
\begin{proof}
Morleyise: $(T^+)^m$ is a Morita extension of $T^m$
\cite[proof of Thm.~5.4]{Tsem17}, so $\cC_{T^m}\to\cC_{(T^+)^m}$ induces
an equivalence of pretopos completions \cite[Cor.~4.9]{Tsem17}; and
$\mathrm{mod}(T)=\Mod(\cC_{T^m})$, $\mathrm{mod}(T^+)=\Mod(\cC_{(T^+)^m})$,
natural transformations between coherent functors on a Morleyisation
being the elementary embeddings. The reduct functor $r\colon
\mathrm{mod}(T^+)\to\mathrm{mod}(T)$ is $\Mod(j)$ for the inclusion
$j\colon\cC_{T^m}\to\cC_{(T^+)^m}$, and $\PP j$ is an equivalence, so
$r$ is an equivalence by Theorem~\ref{thm:cc}. The canonical
expansion $e\colon M\mapsto M^+$ is a functor with $re=\mathrm{id}$
\cite[Thm.~42]{Meadows24}, so $e$ is a quasi-inverse of $r$ and an
equivalence. Now $t^\dagger=\Mod(t^m)\circ e$ with $\Mod(t^m)\colon
\Mod(\cC_{(T^+)^m})\to\Mod(\cC_{S^m})$; if $t^\dagger$ is an equivalence
so is $\Mod(t^m)$, and Theorem~\ref{thm:cc} makes $\PP t^m\colon\PP
\cC_{S^m}\to\PP\cC_{(T^+)^m}\simeq\PP\cC_{T^m}$ an equivalence. Thus
$S^m$ and $T^m$ are Morita equivalent coherent theories
\cite[Thm.~4.7]{Tsem17}, and $S$, $T$ are Morita equivalent
\cite[Thm.~5.4]{Tsem17}. Morita bi-interpretability provides such a
$t^\dagger$ with a quasi-inverse, definable or not.
\end{proof}

The same argument does not apply to the next row of Meadows's table,
\emph{Morita-congruence}, where $s^\dagger t^\dagger(M)\cong M$ is only
required object by object \cite[Def.~44(3)]{Meadows24}: the functor
$t^\dagger$ is then essentially surjective and injective on isomorphism
classes but need not be full or faithful, and conceptual completeness
needs the morphisms (Remark~\ref{rem:bhgap}, Question~\ref{que:groupoid}).
The corollary is to be distinguished from the results relating Morita
equivalence to bi-interpretability in the model-theoretic sense of
\cite[Def.~5.5]{McEl20}, where each theory interprets the models of the
other as definable quotient structures and the two composites are
definably isomorphic to the identity. McEldowney
\cite[Prop.~5.11]{McEl20} proves, for theories each proving that some
sort has two distinct elements, that Morita equivalence \emph{implies}
bi-interpretability, the hypothesis serving to code the coproduct
sorts; he shows that the converse fails for his binary sorts
(Remark~\ref{rem:mcel}) and leaves it open for finite signatures.
Kamsma \cite[Fact~8.6]{Kamsma23} states the two notions to be
equivalent for positively disjoint theories, Morita equivalence being
read there, through \cite{Tsem17}, as having the same classifying
topos, that is as $V_\PP$. Corollary~\ref{cor:meadows} goes in the
direction opposite to McEldowney's, from Meadows's Morita
bi-interpretability, whose interpretations are identity-preserving and
whose quotients and coproducts are supplied by the Morita extensions,
to Morita equivalence; it needs no two-element hypothesis, and its
content is Theorem~\ref{thm:cc}, the Makkai--Reyes theorem, applied to
Meadows's notion.

\section{Conceptual distance as an alternation depth}\label{sec:distance}

\cite{KSLF20} define, for a class $X$ of theories with an equivalence
relation $E$ (there: definitional equivalence) and a symmetric relation
$S$ (there: ``differ by one relation symbol'', $T\subseteq T'$
conservatively with one new symbol), the \emph{step distance} $d(T,T')$
as the least number of $S$-steps in a path $T=T_0,\dots,T_m=T'$ each
step of which is an $E$-step or an $S$-step, and $\infty$ if there is no
path; with these choices it is the \emph{conceptual distance}
\cite[Def.~4.4--4.5]{KSLF20}.

\begin{prop}\label{prop:distance}
Let $E$ and $S$ be as above and let $E\vee S$ be the join in the lattice
of equivalence relations on $X$, the transitive closure of $E\cup S\cup
S^{-1}$. Then $d(T,T')<\infty$ iff $T\,(E\vee S)\,T'$, and $d(T,T')$ is
the alternation depth of the pair: the least $k$ such that $T$ and $T'$
are connected by a chain $E\,S\,E\,S\cdots$ with $k$ letters $S$. In
particular the step distance is bounded on $X$ iff the join $E\vee S$ is
reached by a uniformly bounded alternation, and Theorem~4.7 of
\cite{KSLF20} (every $n\in\mathbb N\cup\{\infty\}$ is a conceptual
distance) states that for conceptual distance some pairs are not related
by $E\vee S$, and that on the related pairs the alternation depth is
finite at each pair but not uniformly bounded.
\end{prop}
\begin{proof}
A path with $k$ letters $S$ is an alternating chain with $k$ $S$-steps
after merging consecutive $E$-steps ($E$ is transitive) and inserting
identity $E$-steps; the two descriptions of the join agree.
\end{proof}

This is the situation of \cite[Thm.~2.3(iii)]{Cutpoints} for two
closure operators: the join of the two classes is the class of the join
iff the alternation is finite at every point, and finiteness at every
point does not give a uniform bound, as in the three-dimensional
template example there. Conceptual distance is the alternation depth
for the pair (definitional equivalence, one-symbol extension), and the
theories of \cite[Thm.~4.7]{KSLF20} with distance $n$ are the analogue
of the sets at which $n$ rounds are needed.

\section{Duality: which comparison levels are recovered from the models}\label{sec:duality}

Maruyama \cite{Maruyama16} asks when a duality between a syntactic side
and a semantic side emerges, changes and breaks. The results above give
an answer for the comparison of theories, in the following form. Say
that a comparison level $W\in\Same(\Coh)$ is \emph{dualisable through}
a $2$-functor $M\colon\Coh^{op}\to\mathcal K$ if $W=V_M$, and that the
duality \emph{holds} at that level if moreover $W=V_R$ for a reflector
$R$ whose image is recovered from $M$ (the localisation at $W$ is
reflective and the reflection is determined by the models).

\begin{thm}\label{thm:duality}
The Morita sameness $V_\PP$ is dualisable through the category of
models, $V_\PP=V_{\Mod}$ (Theorem~\ref{thm:cc}), and through the
ultracategory of models, $V_\PP=V_{\Mod_{\mathrm{ult}}}$; in the second
case the pretopos is recovered as the category of ultrafunctors
$\Mod(\cC)\to\Set$ (Makkai's strong conceptual completeness
\cite{Makkai88}, \cite[Cor.~11]{LurieNotes}), so the duality holds.
\end{thm}
\begin{proof}
The first statement is Theorem~\ref{thm:cc}. For the second,
$\Mod_{\mathrm{ult}}(f)$ is an ultrafunctor, its ultrastructure maps
being invertible, so if its underlying functor is an equivalence a
quasi-inverse carries the transported ultrastructure and
$\Mod_{\mathrm{ult}}(f)$ is an equivalence of ultracategories; thus
$V_{\Mod_{\mathrm{ult}}}=V_{\Mod}$, and $\PP\cC\simeq\mathrm{Fun}_{\mathrm{ult}}(\Mod(\cC),\Set)$
is \cite{Makkai88}.
\end{proof}

The object-level categorical equivalence of \cite{BH16} is the
situation where the semantic side is kept but the morphisms connecting
the syntactic objects are forgotten; the duality then ``breaks'' in the
sense that equivalent model categories no longer come from an
interpretation. Adding structure on the semantic side (ultraproducts, or
the topology of Awodey--Forssell \cite{AF13} on groupoids of models)
restores it.

\begin{que}\label{que:groupoid}
Let $\Mod_{\mathrm{gpd}}(\cC)$ be the groupoid of models and
isomorphisms, the core of $\Mod(\cC)$ without the topology of
\cite{AF13}. Is $V_{\Mod_{\mathrm{gpd}}}=V_\PP$, i.e.\ is a coherent
functor of pretoposes inducing an equivalence of groupoids of models an
equivalence? For $\omega$-categorical theories the automorphism group of
the countable model determines the theory up to bi-interpretability
(Ahlbrandt--Ziegler \cite{AZ86}); in general the topological groupoid of
\cite{AF13} is needed to recover the classifying topos, and Makkai's
\cite{Makkai93} studies which structure on the groupoid of models makes
definability recoverable. Hudetz's Theorem~2 \cite{Hudetz19} is a
statement at this level, and Meadows's Morita-congruence
\cite[Def.~44]{Meadows24}, which constrains only the objects, lies
below it.
\end{que}

\section{Open questions}

\begin{enumerate}
\item Whether three steps suffice for every coherent theory in which every
definable relation between distinct sorts is complemented
(Theorem~\ref{thm:boolean} and the remark after it), and the least
number of steps for a given theory in general.
\item The lattice generated by the five samenesses $V_{\mathrm{prod}}$,
$V_{\mathrm{sub}}$, $V_\FF$, $V_\EE$ and renaming in $\Same$ of the
category of theories with signatures (rather than of $\Coh$, where the
cartesian ones are trivial): which joins are reached in one round, and
the position of definitional equivalence, bi-interpretability
\cite{FV25} and the intermediate notions of \cite{Meadows24}, whose
remaining open arrows (strict bi-interpretability against definitional
equivalence, iso-congruence against strict bi-interpretability,
Problems~21 and 23 of \cite{Meadows24}) are questions about this
lattice.
\item Translations between logics (double negation, Girard, Baaz) as
closure operators on theories, after \cite[\S4.10]{Maruyama16}, and their
joins with the Morita completion.
\end{enumerate}

\section*{Declarations}
The author has no competing interests to declare.

\begin{thebibliography}{99}
\bibitem{AZ86} G.~Ahlbrandt, M.~Ziegler, \emph{Quasi finitely axiomatizable
totally categorical theories}, Ann. Pure Appl. Logic 30 (1986), 63--82.
\bibitem{AF13} S.~Awodey, H.~Forssell, \emph{First-order logical duality},
Ann. Pure Appl. Logic 164 (2013), 319--348.
\bibitem{BH16} T.~W.~Barrett, H.~Halvorson, \emph{Morita equivalence},
Rev. Symb. Log. 9 (2016), 556--582.
\bibitem{CV98} A.~Carboni, E.~M.~Vitale, \emph{Regular and exact
completions}, J. Pure Appl. Algebra 125 (1998), 79--116.
\bibitem{FV25} H.~Friedman, A.~Visser, \emph{When bi-interpretability
implies synonymy}, Rev. Symb. Log. 18 (2025).
\bibitem{HT17} H.~Halvorson, D.~Tsementzis, \emph{Categories of
scientific theories}, in: E.~Landry (ed.), Categories for the Working
Philosopher, Oxford University Press, 2017, 402--429.
\bibitem{Hudetz19} L.~Hudetz, \emph{Definable categorical equivalence},
Philos. Sci. 86 (2019), 47--75.
\bibitem{Elephant} P.~T.~Johnstone, \emph{Sketches of an Elephant: A Topos
Theory Compendium}, Oxford Univ. Press, 2002.
\bibitem{Kamsma23} M.~Kamsma, \emph{Type space functors and
interpretations in positive logic}, Arch. Math. Logic 62 (2023), 1--28.
\bibitem{KSLF20} M.~Khaled, G.~Sz\'ekely, K.~Lefever, M.~Friend,
\emph{Distances between formal theories}, Rev. Symb. Log. 13 (2020),
633--654.
\bibitem{LV01} S.~Lack, E.~M.~Vitale, \emph{When do completion processes
give rise to extensive categories?}, J. Pure Appl. Algebra 159 (2001),
203--230.
\bibitem{LurieNotes} J.~Lurie, \emph{Categorical logic (Math 278X)},
lecture notes, Harvard University, 2018, Lectures 27X--29X.
\bibitem{Makkai88} M.~Makkai, \emph{Strong conceptual completeness for
first-order logic}, Ann. Pure Appl. Logic 40 (1988), 167--215.
\bibitem{Makkai93} M.~Makkai, \emph{Duality and definability in first
order logic}, Mem. Amer. Math. Soc. 503 (1993).
\bibitem{MR77} M.~Makkai, G.~E.~Reyes, \emph{First Order Categorical
Logic}, Lecture Notes in Math. 611, Springer, 1977.
\bibitem{Maruyama16} Y.~Maruyama, \emph{Meaning and Duality: From
Categorical Logic to Quantum Physics}, D.Phil. thesis, University of
Oxford, 2016.
\bibitem{McEl20} P.~A.~McEldowney, \emph{On Morita equivalence and
interpretability}, Rev. Symb. Log. 13 (2020), 388--415.
\bibitem{Meadows24} T.~Meadows, \emph{Beyond linguistic interpretation in
theory comparison}, Rev. Symb. Log. 17 (2024), no.~3, 819--859.
\bibitem{Same} R.~Shibusawa, \emph{The sameness lattice of a category:
weak-equivalence structures, Galois functoriality, and comparison intervals},
preprint (2026), archived at DOI 10.5281/zenodo.22823772.
\bibitem{Fijoin} R.~Shibusawa, \emph{Joins of localisations in the lattice of
two-out-of-three classes: the finite-iteration criterion for closure
operators on exact categories, with applications to toposes, torsion
theories and fuzzy sets}, preprint (2026), archived at DOI
10.5281/zenodo.23691896.
\bibitem{Cutpoints} R.~Shibusawa, \emph{Abstract domains and cut points in
the lattice of two-out-of-three classes: joins of closure operators on a
lattice and of cut-point topologies on a control-flow graph}, preprint
(2026), archived at DOI 10.5281/zenodo.23966743.
\bibitem{Tsem17} D.~Tsementzis, \emph{A syntactic characterization of
Morita equivalence}, J. Symbolic Logic 82 (2017), 1181--1198.
\end{thebibliography}

\end{document}
